
\section{引言}

\subsection{问题背景}

随着人口老龄化加剧和生活方式改变，慢性非传染性疾病已成为全球健康的主要挑战。根据世界卫生组织统计，慢性病患者数量不断增加，其医疗负担和社会成本巨大。高血压、糖尿病、心脏病等慢性疾病需要长期管理和监测，但传统的被动式、间断式管理模式难以满足患者的实际需求。

当前，我国基层医疗机构普遍存在以下困境：患者健康数据分散，无统一管理平台；医患信息沟通不畅，患者依从性低；健康风险评估手段落后，难以实现早期预警；医疗资源配置不合理，管理效率低下。这些问题严重制约了慢性病防控的效果。

为此，迫切需要建立一套现代化的慢性健康管理系统，通过信息技术手段实现患者数据的集中管理、科学分析、实时监测和智能决策支持，提升管理效率，改善患者预后。

\subsection{问题重述}

本文针对慢性健康管理中的关键问题，设计开发了一套基于Python的综合管理系统，主要包括以下四个方面：

\textbf{问题1：患者数据的标准化、清洗与管理。} 如何对来自不同医疗机构、不同检查设备的异质性患者健康数据进行有效的数据清洗、标准化处理，建立统一的数据存储和管理体系，保证数据质量和可用性？

\textbf{问题2：患者健康状态的综合评估与风险预测。} 如何基于多维度的临床指标建立科学的健康评估模型，采用机器学习等先进技术进行疾病风险预测和预警，提高预测准确性和临床实用性？

\textbf{问题3：个性化健康管理方案的智能推荐。} 如何根据患者的个体特征、病情严重程度和生活方式等因素，制定和推荐个性化的管理方案和生活干预建议，提高患者的健康管理依从性？

\textbf{问题4：医患交互与实时健康监测的可视化呈现。} 如何设计用户友好的交互界面，通过丰富的数据可视化手段实现患者数据的实时展示、交互式分析和移动端监测，促进医患沟通和患者自我管理？

\subsection{研究意义}

本研究对完善我国慢性病防控体系、提升基层医疗服务能力具有重要意义，也为患者自我健康管理提供了科学工具，有利于改善国民健康水平。
